\chapter{Conclusiones y trabajo futuro}\label{chap:conclusiones}

Los capítulos anteriores describieron qué hace el PII sobre el levitador ECP-730 en configuración atractiva y por qué lo hace. Este capítulo reúne lo que esos resultados permiten afirmar: responde a la pregunta de investigación planteada en el capítulo~\ref{chap:introduccion}, evalúa la hipótesis, sitúa el trabajo frente a sus antecedentes y propone cómo continuarlo.

\section{Conclusiones}\label{sec:conclusiones}

Los cuatro objetivos específicos de la sección~\ref{sec:objetivos} se cumplieron en el orden en que se plantearon. El análisis de la dinámica (capítulo~\ref{chap:modelo}) mostró que, en la configuración atractiva, cada posición admite un equilibrio y que todos son inestables, y que el atascamiento-deslizamiento puede representarse con la fricción de Karnopp, y la zona muerta, no centrada en el origen, con una retención de voltaje que la simplifica. Ese modelo sirvió de base para el diseño sobre el modelo linealizado (capítulo~\ref{chap:controlador}), cuyo controlador estabiliza el lazo con márgenes holgados y conserva la estabilidad en todo el intervalo de pruebas. La simulación no lineal sin atascamiento-deslizamiento (capítulo~\ref{chap:desempeno}) confirmó que el diseño lineal describe bien a la planta no lineal, de modo que las diferencias que aparecen después no se deben a la linealización. La simulación con atascamiento-deslizamiento, por último, separó lo que el controlador logra, que es la estabilización, de lo que no puede evitar, que es el ciclo límite. Cada objetivo aporta así un eslabón de una misma cadena: sin el modelo no hay diseño, sin el diseño lineal no hay comparación con la planta no lineal, y sin esta no es posible atribuir el ciclo límite a su causa. Esa cadena sostiene la evaluación de la hipótesis que sigue y explica por qué los resultados del capítulo~\ref{chap:doble_integral} pueden leerse como consecuencia del modelo y no como un artefacto de la simulación.

Se diseñó un controlador PII lineal que estabiliza la posición del SLM ECP-730, inestable en lazo abierto, en presencia de zona muerta y fricción estática, sin compensar explícitamente ninguna de las dos. La estabilización se logró con un diseño sobre el modelo linealizado (capítulo~\ref{chap:controlador}) y se verificó por simulación sobre el modelo no lineal con atascamiento-deslizamiento (capítulo~\ref{chap:desempeno}). El precio es un ciclo límite permanente: en estado estacionario la convergencia asintótica se sustituye por una oscilación cuya amplitud es del orden de centésimas de centímetro, con el error medio prácticamente nulo. La tesis responde así que sí es posible estabilizar el sistema con un controlador lineal, pero no sin ciclo límite.

\paragraph{La hipótesis.}
La primera parte de la hipótesis se confirma: el PII estabiliza la posición del imán dentro del intervalo de operación que permite el actuador, tanto en regulación como en seguimiento. La segunda parte no se sostiene de forma general, aunque tampoco se refuta por completo. En las familias PI y PII evaluadas (figura~\ref{fig:r15_familias}), ningún PII aceptado alcanza un ciclo límite menor que el mejor PI, con $0.026$~cm frente a $0.013$~cm, o frente a $0.019$~cm si se excluye también todo recorte de la orden del controlador; el orden se conserva al refinar la integración. Al ampliar la libertad de sintonización de ambos controladores a la vez —su ganancia, su cero de baja frecuencia y su red de adelanto, con un barrido conjunto de 450 combinaciones para cada uno (sección~\ref{sec:pii_v2})— las fronteras de mejor desempeño alcanzable se cruzan cerca de un ciclo límite de $0.008$~cm. Por debajo de ese valor, en esas fronteras optimizadas, el PI iguala o supera al PII en ciclo límite y en atascamiento; por encima, con ganancia baja y el cero de baja frecuencia desplazado lejos de su valor de diseño, ocurre lo contrario: el PII logra atascamientos más cortos que el PI con el mismo ciclo, y en el extremo de menor atascamiento de todo el trabajo lo supera en las dos métricas a la vez. Los controladores de la comparación directa (tabla~\ref{tab:pi_pii}, $0.031$ y $0.033$~cm) no están en el régimen de ciclo pequeño de esas fronteras, sino muy por encima del cruce; ahí el PI también iguala o supera ligeramente al PII, pero por tratarse de una sintonización particular, no optimizada, y esa coincidencia de sentido no debe confundirse con pertenecer al tramo de las fronteras donde domina el PI. La doble acción integral no carece, por tanto, de toda ventaja, pero esa ventaja no aparece en la sintonización nominal de esta tesis ni es, en general, la que anticipaba la hipótesis. No constituye una demostración general sobre PI frente a PII ni una validación experimental.

El análisis del ciclo límite (sección~\ref{sec:analisis_ciclo}) explica por qué, en el régimen habitual, el PII no aventaja al PI. La fricción estática equivale a una banda de voltaje alrededor del equilibrio que la señal de control debe atravesar para que el imán se mueva. La acción integral es la que obliga a cruzar esa banda una y otra vez, y la segunda integración no cambia ese mecanismo por sí sola: en atascamientos de unos $0.35$~s el término integral simple aporta la mayor parte del cambio de voltaje, y la amplitud del ciclo disminuye al aumentar la ganancia de alta frecuencia del controlador, no al añadir un integrador. La derivación de primeros principios de la sección~\ref{sec:fuerza_primeros_principios} y la ecuación de rigidez general~\eqref{eq:rigidez_general} confirman este mecanismo de forma analítica. La segunda integración conserva su ventaja teórica, que es la eliminación del error en rampa, pero esa ventaja queda por debajo del error que impone el propio ciclo límite (tabla~\ref{tab:pi_pii}) en el régimen de ganancia habitual.

\paragraph{Hallazgos adicionales.}
Tres resultados rebasan la hipótesis y tienen valor propio. En primer lugar, el límite del actuador restringe el intervalo de operación: con $u\in[0;\,3.5]$~V el imán puede sostenerse en $-4$~cm, pero no subir desde ahí, y si baja ya no puede recuperarse; los cambios de referencia solo son viables a partir de unos $-3.6$~cm (sección~\ref{sec:intervalo}). Por ello las pruebas se acotaron al intervalo de $-3$ a $-2$~cm, y esa restricción es de la planta, no del controlador. En segundo lugar, el ciclo límite proviene de las no linealidades no suaves y no de la fuerza electromagnética: el modelo linealizado con zona muerta y fricción reproduce un ciclo casi idéntico ($0.028$ frente a $0.033$~cm) al del modelo no lineal completo. En tercer lugar, el mismo PII, diseñado para la configuración inestable, también estabiliza la configuración estable, lo que indica que el diseño no depende de la inestabilidad del equilibrio.

\section{Relación con los antecedentes}

Hernández Alcántara et al.~\cite{alcantara} observaron en la configuración repulsiva del ECP-730 que la acción integral interactúa con la fricción y produce oscilaciones de tipo ciclo límite, y las eliminaron con control conmutado. Este trabajo encuentra el mismo mecanismo en la configuración atractiva y muestra que, cuando se prescinde de la conmutación, un controlador lineal basta para estabilizar el sistema a costa de ese ciclo. La expectativa que dejó el trabajo de Pérez-Gómez~\cite{perez}, según la cual la doble acción integral reduce los efectos de la zona muerta en un motor de CD, no se confirma para el SLM en el régimen de sintonización habitual de este trabajo, aunque sí aparece, acotada, en un régimen de ganancia baja (sección~\ref{sec:pii_v2}). Las plantas difieren, un motor y un levitador inestable, de modo que el resultado delimita el alcance de aquella hipótesis sin invalidar el trabajo que la motivó. Queda abierta la pregunta de si la ventaja de la segunda integración se vuelve dominante en condiciones distintas, por ejemplo con referencias de pendiente mucho mayor o con otros criterios de diseño que acepten el sobrepaso que ese régimen exige.

\section{Trabajo futuro}

La línea de mayor prioridad es la validación experimental, pues todos los resultados provienen de simulación. Implementar el PII en el dispositivo real exige verificar los parámetros del modelo y, si difieren, reidentificarlos mediante ensayos en lazo abierto. En particular, la zona muerta se supuso como hipótesis de modelado y su identificación, por ejemplo con mínimos cuadrados recursivos, permitiría después compensarla de forma activa. Un observador de velocidad extendería la metodología de Hernández Alcántara~\cite{diana_tesis} a la configuración inestable y haría viables las estrategias que requieren la velocidad, como la compensación de la fricción.

La segunda línea ataca directamente el ciclo límite, que es el factor que domina el desempeño. Se propone explorar estructuras que actúen sobre él: la compensación de la fricción estática, el control conmutado entre un PI y un controlador PD, analizado por Pérez-Gómez~\cite{perez} para el motor de CD, y el \emph{dither}. Conviene además rediseñar el controlador en el centro del intervalo de operación, hacia $-2.5$~cm, donde el margen de fase del diseño actual se reduce a $51^\circ$. Para ampliar el intervalo de operación pueden emplearse la ganancia programada, que ajuste el controlador a cada posición, y el control robusto con QFT o $H_\infty$, que acote la incertidumbre paramétrica y cuantifique los márgenes en un equilibrio inestable.

Por último, queda por comparar modelos alternativos de la fuerza electromagnética y de la fricción, como LuGre, Stribeck o Dahl, para saber qué tan sensibles son las conclusiones a las elecciones de modelado del capítulo~\ref{chap:modelo}.
